\documentclass[10pt]{article}
\usepackage[T1]{fontenc}
\usepackage[utf8]{inputenc}
\usepackage[french]{babel}
\usepackage{amsmath,amssymb,mathrsfs,mathpazo}
\usepackage{geometry}
\geometry{a4paper,margin=2cm}
\usepackage{graphicx,array,tabularx,tabularray,colortbl,multirow,multicol,booktabs,enumitem,pifont,ulem,fancybox,diagbox,cancel}
\usepackage{pstricks,pstricks-add,pst-plot,pst-tree,pst-node,pst-text,pst-func,pst-math,pst-bspline,pst-3dplot,pst-eucl,pst-solides3d}
\let\pscalculate\undefined
\usepackage{pst-bezier}
\usepackage{tikz,pgfplots,tkz-tab}
\usetikzlibrary{arrows,arrows.meta,patterns,calc,positioning,trees,decorations.markings,shapes.geometric,shapes.misc,shapes.arrows,intersections,decorations.pathmorphing,decorations.pathreplacing,backgrounds,shadows,babel}
\pgfplotsset{compat=1.18}
\usepgfplotslibrary{fillbetween}
\usepackage[np]{numprint}
\usepackage{listings,makecell,pas-tableur,twoopt}
\let\barre\undefined
\usepackage[pstricks]{tablvar}
\definecolor{darkgreen}{RGB}{0,100,0}
\providecommand{\red}{\color{red}}
\providecommand{\blue}{\color{blue}}
\providecommand{\green}{\color{green}}
\providecommand{\black}{\color{black}}
\providecommand{\white}{\color{white}}
\providecommand{\gray}{\color{gray}}
\providecommand{\cyan}{\color{cyan}}
\providecommand{\magenta}{\color{magenta}}
\providecommand{\yellow}{\color{yellow}}
\providecommand{\violet}{\color{violet}}
\providecommand{\orange}{\color{orange}}
\providecommand{\euro}{\texteuro}
\providecommand{\EUR}{\texteuro}
\providecommand{\ieme}{\textsuperscript{e}}
\providecommand{\degres}{\ensuremath{^\circ}}
\providecommand{\degree}{\ensuremath{^\circ}}
\providecommand{\vv}[1]{\overrightarrow{#1}}
\providecommand{\Coord}[1]{#1}
\providecommand{\decosix}{}
\providecommand{\starredbullet}{\ensuremath{\bullet}}
\pagestyle{empty}
\setlength{\parindent}{0pt}
\usepackage{adjustbox}\DeclareUnicodeCharacter{2236}{:}\DeclareUnicodeCharacter{2190}{\ensuremath{\leftarrow}}
\providecommand{\R}{\mathbb{R}}
\providecommand{\N}{\mathbb{N}}
\providecommand{\Z}{\mathbb{Z}}
\providecommand{\Q}{\mathbb{Q}}
\providecommand{\D}{\mathbb{D}}
\providecommand{\C}{\mathbb{C}}
\providecommand{\e}{{\mathrm{e}}}
\providecommand{\vect}[1]{\overrightarrow{#1}}
\providecommand{\vectt}[1]{\overrightarrow{\mathrm{#1}}}
\providecommand{\vvt}[1]{\overrightarrow{\mathrm{#1}}}
\providecommand{\np}[2][]{{#2~#1}}
\providecommand{\vv}[1]{{\overrightarrow{#1}}}
\renewcommand{\d}{{\mathrm{d}}}
\providecommand{\degres}{{^\circ}}
\providecommand{\degre}{{^\circ}}
\providecommand{\blue}{\color{blue}}
\providecommand{\red}{\color{red}}
\providecommand{\enskip}{\quad}
\providecommand{\llbracket}{\lbrack\!\lbrack}
\providecommand{\rrbracket}{\rbrack\!\rbrack}
\providecommand{\euro}{{\text{€}}}
\providecommand{\no}{{\text{n°}}}
\providecommand{\interval}[2]{{[#1;#2]}}
\providecommand{\intervFO}[2]{{[#1;#2[}}
\providecommand{\dotfill}{\ldots}

\providecommand{\ds}{\displaystyle}
\providecommand{\textding}[1]{#1}
\providecommand{\Oij}{\(\mathrm{O};\vec{\imath},\vec{\jmath}\)}
\providecommand{\Oijk}{\(\mathrm{O};\vec{\imath},\vec{\jmath},\vec{k}\)}
\providecommand{\Ouv}{\(\mathrm{O};\vec{u},\vec{v}\)}
\providecommand{\intervFF}[2]{\ensuremath{[#1;#2]}}
\renewcommand{\euro}{\ifmmode\text{\texteuro}\else\texteuro\fi}
\DeclareUnicodeCharacter{00B0}{\ifmmode^{\circ}\else\textdegree\fi}
\DeclareUnicodeCharacter{20AC}{\ifmmode\text{\texteuro}\else\texteuro\fi}
% Énoncés : Éducation nationale. Transcription et corrigés : APMEP, Denis Vergès et contributeurs. Découpage et classement : CoopMaths, Sébastien Lozano, Éric Elter et contributeurs. Adaptation de la présentation : Maths Explorer.
% CC BY-SA : https://coopmaths.fr/www/about/
% https://coopmaths.fr/alea/static/bac/2025/tex/bac_2025_06_sujet1_asie_3.tex
\begin{document}
\noindent\textit{APMEP / CoopMaths. CC BY-SA. Présentation adaptée par Maths Explorer.}\par\medskip
Un patient doit prendre toutes les heures une dose de $2$ mL d'un médicament.

On introduit la suite $\left(u_n\right)$
telle que le terme $u_n$ représente la quantité de médicament, exprimée en mL présente dans l'organisme immédiatement après $n$ prises de médicament.

On a $u_1 = 2$ et 

\begin{center}pour tout entier naturel $n$ strictement positif : $u_{n+1} = 2 + 0,8u_n$.\end{center}

\subsection*{Partie A}

En utilisant ce modèle, un médecin cherche à savoir à partir de combien de prises du médicament la quantité présente dans l'organisme du patient est strictement supérieure à 9 mL.

\begin{enumerate}
\item Calculer la valeur $u_2$.
\item Montrer par récurrence que :
\[u_n = 10 - 8 \times 0{,}8^{n-1} \,\text{pour tout entier naturel} \:n \:\text{strictement
positif.}\]

\item Déterminer $\lim\limits_{n \to +\infty} u_n$ et donner une interprétation de ce résultat dans le contexte de l'exercice.
\item Soit $N$ un entier naturel strictement positif, l'inéquation $u_N \geqslant 10$ admet-elle des solutions ? 

Interpréter le résultat de cette question dans le contexte de l'exercice.
\item Déterminer à partir de combien de prises de médicament la quantité de médicament présente dans l'organisme du patient est strictement supérieure à $9$~mL. Justifier votre démarche.
\end{enumerate}

\subsection*{Partie B}

En utilisant la même modélisation, le médecin s'intéresse à la quantité moyenne de médicament présente dans l'organisme du malade au cours du temps.

On définit pour cela la suite $\left(S_n\right)$ définie pour tout entier naturel $n$ strictement positif par

\[S_n = \frac{u_1 + u_2 + \dots + u_n}{n}.\]

On admet que la suite $\left(S_n\right)$ est croissante.

\medskip

\begin{enumerate}
\item Calculer $S_2$.
\item Montrer que pour tout entier naturel $n$ strictement positif,

\[u_1 + u_2 + \dots + u_n = 10n - 40 + 40 \times 0{,}8^n.\]

\item Calculer $\lim\limits_{n \to +\infty} S_n$.
\item On donne la fonction mystere suivante, écrite en langage Python :

\begin{center}
\begin{tabular}{c|l|}\cline{2-2}
1&def mystere(k):\\
2&\quad n = 1\\
3&\quad s = 2\\
4&\quad while s < k:\\
5&\quad\qquad n = n + 1\\
6&\quad\qquad s = 10 - 40/n + (40*0.8**n)/n\\
7&\quad return n\\ \cline{2-2}
\end{tabular}
\end{center}

Dans le contexte de l'énoncé, que représente la valeur renvoyée par la saisie \texttt{mystere(9)} ? 
\item Justifier que cette valeur est strictement supérieure à $10$.
\end{enumerate}

\end{document}
